\documentclass[10pt]{article}
\usepackage[margin=0.9in]{geometry}
\usepackage{enumitem}
\usepackage{booktabs}
\usepackage{hyperref}
\setlist{nosep,leftmargin=*}
\setlength{\parskip}{4pt}
\setlength{\parindent}{0pt}

\title{\vspace{-1.5em}Matrix-Multiply Accelerator Integration in a RISC-V SoC\vspace{-0.5em}}
\author{Progress Report}
\date{\vspace{-1em}\today}

\begin{document}
\maketitle
\vspace{-1em}

\section*{Objective}
Attach a hardware matrix-multiply accelerator ($C = A \times B$, 32-bit signed
operands and accumulators) to an existing RISC-V SoC. The CPU configures it
through memory-mapped registers, and it accesses main memory on its own to
read $A$ and $B$ and write $C$.

\section*{Work Completed}
We designed a matrix-multiply accelerator in Verilog and connected it to
the SoC's on-chip bus (AXI). The accelerator appears to the processor as a
small set of memory-mapped registers. Software writes in where matrices
$A$, $B$ and $C$ are stored in memory and their dimensions, then sets a
start bit. The accelerator then works on its own: it reads $A$ and $B$
from main memory, computes the product on a $4\times4$ array of
multiply-accumulate units, writes $C$ back to memory, and sets a
``done'' flag. To make this possible, the SoC's address decoder was
extended so the accelerator's registers have their own address range, and
the bus arbiter was configured to let the accelerator share main memory
with the processor.

To verify the design, we wrote a bare-metal C program that runs on the
RISC-V processor. It configures the accelerator, waits for it to finish
and compares the result with a software-computed reference. The whole
system (processor, bus, memory and accelerator) was run in a cycle-accurate
simulation built with Verilator and SystemC. Getting this working required
fixes to how bus responses are routed back to the accelerator and to the
simulation environment. The simulated system now computes a
$2\times3$ by $3\times2$ matrix product correctly and the test passes.

\section*{Work To Be Done}
\begin{itemize}
  \item Support larger matrices (tiling and burst transfers beyond the
  current $4\times4$ limit).
  \item Test on FPGA hardware.
  \item Measure and report performance: cycle counts, speedup over a
  software-only matrix multiply on the CPU, and resource and energy
  efficiency.
\end{itemize}

\end{document}
